%------------------------------------------------------------------------------%
\documentclass[fleqn,utf8,aspectratio=169,12pt]{beamer}

\usepackage{gaal}

\title{Determinantes e Sistemas Lineares}

%------------------------------------------------------------------------------%
\begin{document}

\addtitle

%------------------------------------------------------------------------------%
\section{Determinantes e Sistemas Lineares}

%------------------------------------------------------------------------------%
\begin{frame}{Determinantes e Sistemas Lineares}
  \onslide<+->{O determinante de $A$ }
  \\
  \onslide<+->{permite responder questões sobre o sistema
  \(
    A \mathbf{x} = \mathbf{b}
  \)}
  \begin{itemize}[<+->]
      \item O sistema possui solução?
      \item A solução é única?
      \item Existem infinitas soluções?
      \item O sistema homogêneo possui soluções não triviais?
  \end{itemize}
\end{frame}

%------------------------------------------------------------------------------%
\begin{frame}{Relação Fundamental}
  \onslide<+->{Relação fundamental
  \[
    \det(A)\neq0
    \qquad\Leftrightarrow\qquad
    A\text{ é invertível}
  \]}

  \onslide<+->{Se
  \(
    \;\det(A) \neq 0\;
  \)}
  \begin{align*}
    \onslide<+->{A \mathbf{x}        & = \mathbf{b}        \\}
    \onslide<+->{A^{-1} A \mathbf{x} & = A^{-1} \mathbf{b} \\}
    \onslide<+->{\mathbf{x}          & = A^{-1} \mathbf{b}
    \qquad\text{ solução única}}
  \end{align*}
\end{frame}

%------------------------------------------------------------------------------%
\begin{frame}{Classificação de sistemas}
  \onslide<+->{Um sistema linear quadrado
  \(
    \;A \mathbf{x} = \mathbf{b}\;
  \)}
  \onslide<+->{pode ser}
  \begin{itemize}[<+->]
      \item possível e determinado
      \item possível e indeterminado
      \item impossível
  \end{itemize}

  \bigskip

  \onslide<+->{Se
  \(
    \;\det(A) \neq 0\;
  \)}
  \onslide<+->{o sistema é necessariamente \emph{possível e determinado}}

  \onslide<+->{Se
  \(
    \;\det(A) = 0\;
  \)}
  \onslide<+->{o sistema pode ser \emph{possível e indeterminado} ou \emph{impossível}}
\end{frame}

%------------------------------------------------------------------------------%
%------------------------------------------------------------------------------%
\begin{question}
  Utilize o determinantes para verificar se o sistema possui solução única
  \[
    \begin{cases}
      2x + y = 5 \\
       x - y = 1
    \end{cases}
  \]
\end{question}

%------------------------------------------------------------------------------%
\begin{resolution}{Solução}
  Matriz dos coeficientes
  \[
    A =
    \begin{pmatrix}
      2 & 1  \\
      1 & -1
    \end{pmatrix}
  \]

  \pause
  Determinante
  \[
    \det(A)
    \pause
    \;=\; 2(-1) - 1(1)
    \pause
    \;=\; -3
    \pause
    \;\neq\; 0
  \]
  \pause
  o sistema possui uma única solução
\end{resolution}

%------------------------------------------------------------------------------%
\endinput



%------------------------------------------------------------------------------%
\section{Regra de Cramer}

%------------------------------------------------------------------------------%
\begin{frame}{Regra de Cramer}
  Dado o sistema
  \(
    A \mathbf{x} = \mathbf{b}
  \)

  \pause
  com $A$ é quadrada e
  \(
    \det(A) \neq 0
  \)

  \pause
  \bigskip
  A \emph{Regra de Cramer} fornece a solução do sistema por determinantes
\end{frame}

%------------------------------------------------------------------------------%
\begin{frame}{Regra de Cramer para Sistemas $2 \times 2$}
  \[
    \begin{cases}
      a_{11}\, x + a_{12}\, y = b_1  \\
      a_{21}\, x + a_{22}\, y = b_2
    \end{cases}
  \]
  \pause
  \[
    x = \frac{\det(A_x)}{\det(A)}
    \qquad\qquad
    y = \frac{\det(A_y)}{\det(A)}
  \]

  \pause
  $A_x$ e $A_y$ são obtidas substituindo a coluna correspondente por $\mathbf{b}$
  \pause
  \[
    A_x =
    \begin{pmatrix}
      \emph{b_1} & a_{12} \\
      \emph{b_2} & a_{22}
    \end{pmatrix}
  \pause
  \qquad\qquad
    A_y =
    \begin{pmatrix}
      a_{11} & \emph{b_1} \\
      a_{21} & \emph{b_2}
    \end{pmatrix}
  \]
\end{frame}

%------------------------------------------------------------------------------%
%------------------------------------------------------------------------------%
\begin{question}
  Resolva o sistema pela Regra de Cramer
  \[
    \left\{\begin{array}{rrr}
      2x & +y & = 5 \\
       x & -y & = 1
    \end{array}\right.
  \]

\end{question}

%------------------------------------------------------------------------------%
\begin{resolution}{Solução}
  \onslide<+->{\[
    A =
    \begin{pmatrix}
      2 & 1  \\
      1 & -1
    \end{pmatrix}
    \qquad\qquad
    \mathbf{b} =
    \begin{pmatrix}
      5 \\
      1
    \end{pmatrix}
  \]}
  \begin{align*}
    \onslide<+->{\det(A)}
    \onslide<+->{& = 2(-1) - 1 \times 1 \\}
    \onslide<+->{& = -2 -1 \\}
    \onslide<+->{& = -3}
  \end{align*}
\end{resolution}

%------------------------------------------------------------------------------%
\begin{resolution}{Determinando $x$}
  \[
    A =
    \begin{pmatrix}
      2 & 1  \\
      1 & -1
    \end{pmatrix}
  \qquad\qquad
    \mathbf{b} =
    \begin{pmatrix}
      5 \\
      1
    \end{pmatrix}
  \qquad\qquad
    A_x =
    \begin{pmatrix}
      5 & 1  \\
      1 & -1
    \end{pmatrix}
  \]

  \pause
  \[
    \det(A_x)
    \pause
    \;=\; 5(-1) - 1(1)
    \pause
    \;=\; -6
  \]

  \[
  \pause
    x
    \pause
    \;=\; \frac{\det(A_x)}{\det(A)}
    \pause
    \;=\; \frac{-6}{-3}
    \pause
    \;=\; 2
  \]
\end{resolution}

%------------------------------------------------------------------------------%
\begin{resolution}{Determinando $x$}
  \[
    A =
    \begin{pmatrix}
      2 & 1  \\
      1 & -1
    \end{pmatrix}
  \qquad\qquad
    \mathbf{b} =
    \begin{pmatrix}
      5 \\
      1
    \end{pmatrix}
  \qquad\qquad
    A_y =
    \begin{pmatrix}
      2 & 5 \\
      1 & 1
    \end{pmatrix}
  \]

  \pause
  \[
    \det(A_y)
    \pause
    \;=\; 2 \times 1 - 5 \times 1
    \pause
    \;=\; -3
  \]

  \[
    y
    \pause
    \;=\; \frac{\det(A_y)}{\det(A)}
    \pause
    \;=\; \frac{-3}{-3}
    \pause
    \;=\; 1
  \]
\end{resolution}

%------------------------------------------------------------------------------%
\endinput


%------------------------------------------------------------------------------%
\begin{frame}{Regra de Cramer para Sistemas $3 \times 3$}
  \[
    A
    \begin{pmatrix}
      x \\
      y \\
      z
    \end{pmatrix}
    =
    \begin{pmatrix}
      b_1 \\
      b_2 \\
      b_3
    \end{pmatrix}
  \]

  \pause
  Se
  \(
  \det(A)\neq0
  \)
  então

  \pause
  \[
    x = \frac{\det(A_x)}{\det(A)}
    \qquad\qquad
    y = \frac{\det(A_y)}{\det(A)}
    \qquad\qquad
    z = \frac{\det(A_z)}{\det(A)}
  \]
\end{frame}

%------------------------------------------------------------------------------%
%------------------------------------------------------------------------------%
\begin{question}
  Resolva pelo Método de Cramer
  \[
    \begin{cases}
      x   +y + z = 6 \\
      2x  -y + z = 3 \\
      x  +2y - z = 2
    \end{cases}
  \]
\end{question}

%------------------------------------------------------------------------------%
\begin{resolution}{Solução}
  \onslide<+->{\[
    A =
    \begin{pmatrix}
      1 & 1  & 1  \\
      2 & -1 & 1  \\
      1 & 2  & -1
    \end{pmatrix}
  \]}

  \onslide<+->{Pela Regra de Sarrus}
  \begin{align*}
    \onslide<+->{\det(A) & =}
    \onslide<+->{    1(-1)(-1)}
    \onslide<+->{  + 1(1)(1) }
    \onslide<+->{  + 1(2)(2)}
    \onslide<+->{  - 1(-1)(1)}
    \onslide<+->{  - 1(2)(-1)}
    \onslide<+->{  - 1(1)(2)}
    \\
    \onslide<+->{& = 1 + 1 + 4 + 1 + 2 - 2 \\}
    \onslide<+->{& = 7}
  \end{align*}
  \onslide<+->{Como
  \(
    \det(A) \neq 0
  \)}
  \onslide<+->{o sistema possui uma única solução}
\end{resolution}

%------------------------------------------------------------------------------%
\begin{resolution}{Cálculo do Determinante de $A_x$}
  \onslide<+->{\[
    A_x =
    \begin{pmatrix}
      6 & 1  & 1  \\
      3 & -1 & 1  \\
      2 & 2  & -1
    \end{pmatrix}
  \]}
  \begin{align*}
    \onslide<+->{\det(A_x)}
    \onslide<+->{& = 6(-1)^{1+1}
      \begin{vmatrix}
        -1 & 1  \\
        2  & -1
      \end{vmatrix}}
    \onslide<+->{  + 1 (-1)^{1+2}
      \begin{vmatrix}
        3 & 1  \\
        2 & -1
      \end{vmatrix}}
    \onslide<+->{  + 1 (-1)^{1+3}
      \begin{vmatrix}
        3 & -1 \\
        2 & 2
      \end{vmatrix}}
    \\
    \onslide<+->{& = 6\big[(-1)(-1)  -1 \times 2\big]}
    \onslide<+->{  -  \big[3(-1) - 1 \times 2\big]}
    \onslide<+->{  +  \big[3\times 2 - (-1)2 \big] \\}
    \onslide<+->{& = -6 + 5 + 8  \\}
    \onslide<+->{& = 7}
  \end{align*}
\end{resolution}

%------------------------------------------------------------------------------%
\begin{resolution}{Cálculo do Determinante de $A_y$}
  \onslide<+->{\[
    A_y =
    \begin{pmatrix}
      1 & 6 & 1  \\
      2 & 3 & 1  \\
      1 & 2 & -1
    \end{pmatrix}
  \]}
  \begin{align*}
    \onslide<+->{\det(A_y)}
    \onslide<+->{& = 1 (-1)^{1+1}
        \begin{vmatrix}
          3 & 1  \\
          2 & -1
        \end{vmatrix}}
    \onslide<+->{  + 6 (-1)^{1+2}
        \begin{vmatrix}
          2 & 1  \\
          1 & -1
        \end{vmatrix}}
    \onslide<+->{  + 1 (-1)^{1+3}
        \begin{vmatrix}
        2&3\\
        1&2
        \end{vmatrix} \\}
    \onslide<+->{& = 1\big[3(-1) - 1 \times 2\big] }
    \onslide<+->{  - 6\big[2(-1) - 1 \times 1\big]}
    \onslide<+->{  +  \big[2 \times 2 - 3 \times 1\big] \\}
    \onslide<+->{& = -5+18+1  \\}
    \onslide<+->{& = 14}
  \end{align*}
\end{resolution}

%------------------------------------------------------------------------------%
\begin{resolution}{Cálculo do Determinante de $A_z$}
  \onslide<+->{\[
    A_z =
    \begin{pmatrix}
      1 & 1  & 6 \\
      2 & -1 & 3 \\
      1 & 2  & 2
    \end{pmatrix}
  \]}
  \begin{align*}
    \onslide<+->{\det(A_z)}
    \onslide<+->{& = 1 (-1)^{1+1}
        \begin{vmatrix}
          -1 & 3 \\
          2  & 2
        \end{vmatrix}}
    \onslide<+->{  + 1 (-1)^{1+2}
        \begin{vmatrix}
          2 & 3 \\
          1 & 2
        \end{vmatrix}}
    \onslide<+->{  + 6 (-1)^{1+3}
        \begin{vmatrix}
          2 & -1 \\
          1 & 2
        \end{vmatrix} \\}
    \onslide<+->{& = 1\big[(-1)2 - 3 \times 2\big] }
    \onslide<+->{  -  \big[2 \times 2 - 3 \times 1\big]}
    \onslide<+->{  + 6\big[2 \times 2 - (-1)1\big] \\}
    \onslide<+->{& = -8 -1 +30 \\}
    \onslide<+->{& = 21}
  \end{align*}
\end{resolution}

%------------------------------------------------------------------------------%
\begin{resolution}{Solução}
\[
  x
  \pause
  \;=\; \frac{\det(A_x)}{\det(A)}
  \pause
  \;=\; \frac{7}{7}
  \pause
  \;=\; 1
\]

\pause
\[
  y
  \pause
  \;=\; \frac{\det(A_y)}{\det(A)}
  \pause
  \;=\; \frac{14}{7}
  \pause
  \;=\; 2
\]

\pause
\[
  z
  \pause
  \;=\; \frac{\det(A_z)}{\det(A)}
  \pause
  \;=\; \frac{21}{7}
  \pause
  \;=\; 3
\]
\end{resolution}

%------------------------------------------------------------------------------%
\endinput

\endinput
%------------------------------------------------------------------------------%
\begin{question}
\end{question}

%------------------------------------------------------------------------------%
\begin{resolution}{Solução}
\end{resolution}

%------------------------------------------------------------------------------%
\endinput



\begin{frame}{Exemplo -- SPI}
Considere
\[
\begin{cases}
x+2y=3\\
2x+4y=6.
\end{cases}
\]
\medskip

As duas equações representam a mesma reta.
\medskip

A segunda equação é o dobro da primeira:
\[
2(x+2y)=2(3).
\]
\medskip

Portanto, existem infinitas soluções.
\[
x+2y=3.
\]
Tomando $y=t$:
\[
x=3-2t.
\]
\medskip

Assim,
\[
\boxed{
(x,y)=(3-2t,t),\quad t\in\mathbb{R}.
}
\]
\end{frame}

\begin{frame}{Exemplo -- SI}
Agora considere
\[
\begin{cases}
x+2y=3\\
2x+4y=7.
\end{cases}
\]
\medskip

A matriz dos coeficientes é a mesma:
\[
A=
\begin{pmatrix}
1&2\\
2&4
\end{pmatrix},
\]
portanto,
\[
\det(A)=0.
\]
\medskip

Entretanto, o sistema é incompatível.
\end{frame}

\begin{frame}{Exemplo -- SI}
Multiplicando a primeira equação por $2$:
\[
2x+4y=6.
\]
\medskip

Mas a segunda equação afirma:
\[
2x+4y=7.
\]
\medskip

Subtraindo:
\[
0=1,
\]
uma contradição.
\medskip

Portanto,
\[
\boxed{\text{o sistema não possui solução}}.
\]
\end{frame}


\endinput
%------------------------------------------------------------------------------%
\begin{question}
\end{question}

%------------------------------------------------------------------------------%
\begin{resolution}{Solução}
\end{resolution}

%------------------------------------------------------------------------------%
\endinput


\begin{frame}{Exemplo -- Sistema homogêneo}
Considere
\[
\begin{cases}
x+y+z=0\\
2x+2y+2z=0\\
x-y+z=0.
\end{cases}
\]
\medskip

A matriz dos coeficientes é
\[
A=
\begin{pmatrix}
1&1&1\\
2&2&2\\
1&-1&1
\end{pmatrix}.
\]
\medskip

A segunda linha é o dobro da primeira. Portanto,
\[
\det(A)=0.
\]
\medskip

Logo, existem soluções não triviais.
\end{frame}

\begin{frame}{Exemplo -- Sistema homogêneo}
As equações independentes são
\[
x+y+z=0
\]
e
\[
x-y+z=0.
\]
\medskip

Subtraindo:
\[
2y=0,
\]
portanto,
\[
y=0.
\]
\medskip

Da primeira equação:
\[
x+z=0.
\]
Logo,
\[
x=-z.
\]
\medskip

Tomando
\[
z=t,
\]
temos
\[
(x,y,z)=(-t,0,t).
\]
\end{frame}



%------------------------------------------------------------------------------%
% \framedgraphic{Gráfico de uma Função Real}{B_01_grafico_funcao_real}{}

%------------------------------------------------------------------------------%
%\section{Lista Mínima}
%
%%------------------------------------------------------------------------------%
%\begin{frame}{Lista Mínima}
%  \begin{enumerate}
%    \item
%  \end{enumerate}
%
%  \vfill
%  Atenção: A prova é baseada no livro, não nas apresentações
%\end{frame}

%------------------------------------------------------------------------------%
\end{document}
%------------------------------------------------------------------------------%
